\documentclass[10pt,a4paper]{amsart}
\usepackage[margin=19mm]{geometry}
\usepackage{amsmath,amssymb,mathtools}
\usepackage[scr=rsfs,cal=euler]{mathalfa}
\usepackage{xfrac}
\usepackage[unicode,hidelinks,bookmarks=false]{hyperref}
\newcommand{\ie}{i.e.\ }
\newcommand{\cf}{cf.\ }
\newcommand{\resp}{respectively\ }
\newcommand{\Subsection}{\subsection}
\providecommand{\nobreakdash}{\nobreak\hbox{-}}
\setlength{\emergencystretch}{2em}
\multlinegap0pt
\usepackage{dsfont}
\usepackage{amssymb}
\usepackage{bm}
\usepackage{mathtools}
\usepackage{booktabs}
\usepackage[shortlabels]{enumitem}
\usepackage{algorithm}\usepackage{algpseudocode}
\usepackage{xcolor}
\usepackage{tikz}
\usepackage{tcolorbox}
\newtheorem{lemma}{Lemma}
\DeclareMathOperator{\Aut}{Aut}
\title{Implementation of Oblivious Transfer over Binary-Input AWGN Channels by Polar Codes \\}
\author{Pin-Hsun Lin, Hadi Aghaee, Christian Deppe, Eduard A. Jorswieck, Holger Boche}
\date{Source: arXiv:2601.10682v2 (CC BY 4.0)}
\begin{document}
\maketitle

\noindent\emph{Source and licence.} Implementation of Oblivious Transfer over Binary-Input AWGN Channels by Polar Codes \\. Version arXiv:2601.10682v2 (CC BY 4.0), distributed by arXiv under the Creative Commons Attribution 4.0 International licence, http://creativecommons.org/licenses/by/4.0/, as recorded in the arXiv licence metadata for this exact version and shown on the abstract page https://arxiv.org/abs/2601.10682v2. Full licence notice: \texttt{licenses/arxiv-2601.10682-CC-BY-4.0.txt}.\par\smallskip

\noindent\textit{Editorial scope.} Counted unit: the article's lemma lem:polar-incidence (source0.tex lines 1757--1763). The article's complete original proof of it is delivered with it (source0.tex lines 3891--3940, one \texttt{proof} environment of 50 lines, which the article places away from the statement). No mathematics is written, repaired or re-derived: the statement and the proof are the article's own lines, in the article's own notation, and no retained line is substituted or re-flowed.  Retained uncounted as the source-local context the proof needs: lines 1771--1779.  Nothing was omitted from any retained span, so the package carries no omission record.  No retained line contains a citation, so the wrapper carries no bibliography and no bibliography entry.  No retained line contains a figure environment, an image-inclusion command or drawing code, so no image is delivered and the package declares no asset dependency.  Editorial formatting only, in addition: an AMS \texttt{amsart} preamble that restates the article's own declarations and macros used by the retained lines, namely the environments \texttt{lemma} on separate counters, together with the standard packages those lines need.  The retained environments print in this document as Lemma 1 in order of appearance, whereas the article prints the same items with the numbering of its own full document.  The complete native source of the article as distributed in the arXiv e-print for this exact version is bundled unchanged as \texttt{sources/192793/source0.tex}.
\begin{align}\label{EQ_P_pi-compact}
  (\mathbf P_\pi)_{u,x}
  \;=\;
  \begin{cases}
    1, & u = \pi(x),\\
    0, & \text{otherwise},
  \end{cases}
  \qquad u,x\in\mathcal X.
\end{align}

\begin{lemma}\label{lem:polar-incidence}
For every $m\ge 1$ and every $x,y\in\mathcal{X}$,
\begin{equation}\label{eq:T-encodes-order-compact}
  \mathbf T_{x,y}
  \;=\;  \mathds{1}\{\,y \le_b x\,\}.
\end{equation}
\end{lemma}

\begin{proof}
Recall that \(\mathcal{X}=\{0,1\}^m\) with bit-wise order
\(x\le_b y \mbox{ iff } x_i\le y_i\) for all \(i\in[m]\), and that by
Lemma~\ref{lem:polar-incidence},
\(
  \mathbf{T}_{x,y} = \mathds{1}\{\,y\le_b x\,\}
\)
for all \(x,y\in\mathcal{X}\). Let \(\pi:\mathcal{X}\to\mathcal{X}\) be a bijection, and let
\(\mathbf{P}_\pi\) be the corresponding permutation matrix, whose
\(x\)-th column is \(e_{\pi(x)}\), i.e., \eqref{EQ_P_pi-compact}. 
Then \((\mathbf{P}_\pi^{\top})_{x,u} = (\mathbf{P}_\pi)_{u,x}\).

For any \(x,y\in\mathcal{X}\), we can derive
\begin{align}
  \bigl(\mathbf{P}_\pi^{\top}\mathbf{T}\,\mathbf{P}_\pi\bigr)_{x,y}
  &= \sum_{u,v}
     (\mathbf{P}_\pi^{\top})_{x,u}\,\mathbf{T}_{u,v}\,(\mathbf{P}_\pi)_{v,y}\notag\\
  &= \sum_{u,v}
     (\mathbf{P}_\pi)_{u,x}\,\mathbf{T}_{u,v}\,(\mathbf{P}_\pi)_{v,y}\notag\\
  &= \sum_{u,v}
     \mathds{1}\{u=\pi(x)\}\,\mathbf{T}_{u,v}\,\mathds{1}\{v=\pi(y)\}\label{EQ_Aut_expansion}\\
  &= \mathbf{T}_{\pi(x),\pi(y)}\label{EQ_Aut_expansion2}\\
  &= \mathds{1}\{\,\pi(y)\le_b \pi(x)\,\}.  \label{eq:PpiTPpi-entry}
\end{align}

Assume \(\mathbf{P}_\pi^{\top}\mathbf{T}\,\mathbf{P}_\pi=\mathbf{T}\), then for all \(x,y\in\mathcal{X}\), we have
\[
  \mathds{1}\{\,\pi(y)\le_b \pi(x)\,\}
  = \bigl(\mathbf{P}_\pi^{\top}\mathbf{T}\,\mathbf{P}_\pi\bigr)_{x,y}
  = \mathbf{T}_{x,y}
  = \mathds{1}\{\,y\le_b x\,\},
\]
where the first equality is from \eqref{eq:PpiTPpi-entry}. Hence, we have $ y\le_b x  \,\mbox{ iff }\,  \pi(y)\le_b \pi(x)  \,\forall\,x,y\in\mathcal{X}.$ 

Conversely, assume  $  x\le_b y \mbox{ iff } \pi(x)\le_b \pi(y),  \,\forall\,x,y\in\mathcal{X}.$
Then for all \(x,y\), we have $  \mathds{1}\{\,\pi(y)\le_b \pi(x)\,\}  = \mathds{1}\{\,y\le_b x\,\}  = \mathbf{T}_{x,y}.$
Comparing with \eqref{eq:PpiTPpi-entry}, we obtain
\((\mathbf{P}_\pi^{\top}\mathbf{T}\,\mathbf{P}_\pi)_{x,y} = \mathbf{T}_{x,y}\)
for all \(x,y\), i.e.,
\(
  \mathbf{P}_\pi^{\top}\mathbf{T}\,\mathbf{P}_\pi = \mathbf{T}.
\)


Finally, the map \(\pi\mapsto \mathbf{P}_\pi\) is one-to-one and onto between the two sets:
each bijection \(\pi\) corresponds to exactly one permutation matrix \(\mathbf{P}_\pi\), and vice versa.
Therefore, the equivalence above gives a one-to-one correspondence between \(\Aut(\mathbf{T})\)
and \(\Aut(\mathcal{X},\le_b)\).

\end{proof}

\end{document}
